\subsection{Common bases: from a matroid representation to physical random tapes}
\label{sec:common-bases}

The common-bases development follows an approximate count all the way to a
fixed oracle program, its bit costs, and the tape positions supplying its
randomness \citep{OAICommonBases}. Three descriptions of the computation
are kept explicit: the mathematical estimator, a finite random experiment,
and the compiled machine. Successive interfaces track the counted objects,
the output law, and execution on a given input through this pipeline.

\paragraph{Encoding the counting problem.}
The mathematical problem is to count the common bases of two equal-rank matroids $M_1,M_2$ on an enumerated ground set $E=\operatorname{Fin}n$. The first representation replaces two simultaneous base conditions by one base condition. On $E\sqcup E$, set
\[
 N=M_1\oplus M_2^*,\qquad
 T(A)=\operatorname{inl}(A)\cup\operatorname{inr}(E\setminus A).
\]
The source proves that $T(A)$ is a base of $N$ exactly when $A$ is a base of both input matroids.\footnote{\href{https://github.com/openai/math/blob/fd4aeeb2ee4fc729c18d98444fed42fd0529eeeb/lean/OAI/Combinatorics/MatroidCounting/CommonBases.lean\#L7156-L7167}{Base correspondence and the rank of the paired matroid}.} Equal ranks give $r(N)=n$, while every transversal $T(A)$ has size $n$. Consequently the deficiency $d(A)=n-r_N(T(A))$ vanishes precisely at the objects being counted. Explicit equivalences between subsets, finite subsets, and transversals transfer both cardinalities and partition sums. If $Z$ is the common-base count, the resulting annealing partition function satisfies
\[
 Z\leq \sum_{A\subseteq E}q^{d(A)}\leq Z+2^nq,
 \qquad 0\leq q\leq1.
\]
The partition function now has a direct relation to the desired count:
zero-deficiency terms contribute exactly $Z$, and every other term is at
most $q$. The doubled ground set also supplies the exchange space. A plain
state selects one element from each pair; a defect state has one empty
pair, one full pair, and singleton choices elsewhere. The walk passes
through paired defects before a hitting construction returns it to plain
states. This separates the space convenient for local moves from the
transversals whose weighted count is estimated.

A greedy routine next computes rank from independence queries. Its
correctness theorem identifies the implemented deficiency and annealing
weights with the mathematical ones. A sampled zero-deficiency witness
serves as the estimator's final guard. Because such a witness is a common
base, this guard enforces the zero-count conclusion on every execution,
including tapes outside the high-probability accuracy event.

\paragraph{From experiments to a bounded program.}
Finite random experiments are trees with uniform choices at their internal
nodes. The bounded implementation gives each draw a bank of rejection
trials. Its fit condition bounds the number of successive draws by $M$
and requires $H\geq\lceil\log_2 a\rceil$ trial bits for each positive draw
arity $a$. With $D$ trials per draw, the observable loss is at most
$M2^{-D}$ for values in $[0,1]$. The width condition ensures that a trial
can represent every outcome; increasing the number of undersized trials
would not repair that defect. A separate reindexing relation permits
permutations of outcome labels and proves expectation preservation.
Consequently, the walk and its implementation may enumerate choices
differently while agreeing in law.

In the expression language, randomness becomes explicit input. Evaluation
takes a value and a random bank and returns an optional result together
with the unused bank. A realization theorem relates this deterministic
execution, bank by bank, to the random tree. For a Metropolis step, the
proof follows the proposal, acceptance calculation, Bernoulli choice, and
lazy branch. Composition with the witness guard then gives the complete
estimator. Law preservation and bankwise execution preservation enter at
different layers of the same construction.

Compilation preserves the encoded output and supplies a cost bound.
The constructed nine-register stack machine implements the expression
language by structural simulation, including its list folds. Source
evaluation returns both a value and a charge; the machine produces that
value within the recorded charge. This makes resource use part of the
simulation statement, ready for the polynomial-bound argument.\footnote{\href{https://github.com/openai/math/blob/fd4aeeb2ee4fc729c18d98444fed42fd0529eeeb/lean/OAI/Combinatorics/MatroidCounting/CommonBases.lean\#L16459-L16620}{Cost-bearing compiler correctness}.}

\paragraph{From semantic preservation to a polynomial cost bound.}
Bounding the charge requires control of intermediate representations. The recursive
predicate $\operatorname{Safe}(O,S,e,x)$, with oracle $O$, width bound
$S$, expression $e$, and input $x$, bounds the encoded widths of the input and
every successful output, and requires safety of the subevaluations actually
used. For a composition it bounds the intermediate result; for a fold it
requires the invariant after each accumulator update. If the body charge
is at most $E$ and these widths are at most $S$, the fold lemma gives
\[
  c_{\mathrm{fold}}\leq |\mathrm{xs}|\,(7+12S+E).
\]
The encoding bounds list length by width, so this estimate can be composed
through structural induction on the fixed expression. The width
certificate turns the cost-bearing simulation into a polynomial charge
bound.\footnote{See
\href{https://github.com/openai/math/blob/fd4aeeb2ee4fc729c18d98444fed42fd0529eeeb/lean/OAI/Combinatorics/MatroidCounting/CommonBases.lean\#L16855-L16904}{\src{Safe} and the fold-charge proof}, and
\href{https://github.com/openai/math/blob/fd4aeeb2ee4fc729c18d98444fed42fd0529eeeb/lean/OAI/Combinatorics/MatroidCounting/CommonBases.lean\#L17079-L17092}{the structural polynomial bound}.}

The source packages these certificates as \src{PolySafe}: a polynomial
$p$ takes an input parameter $S$ to an admissible width $p(S)$.
Polynomial substitution supports composition. A further interface,
\src{AttrSafe}, carries numerical predicates alongside the width bound.
This matters for value-indexed operations: an exponent can be very large
despite a short binary encoding. The exponentiation lemma consequently
assumes that the exponent itself is at most $S$. The estimator's attribute
bounds three parameter fields, and its proof transports those bounds
through exponentiation, range construction, and ratio products.
In \src{attrSafe_compose}, the next width parameter is enlarged to
$g(S)=p(S)+S$; monotonicity of the numerical attribute carries its
information to that larger parameter. Thus a usable complexity interface
preserves both encoded size and the value constraints needed by the next
operation.\footnote{See
\href{https://github.com/openai/math/blob/fd4aeeb2ee4fc729c18d98444fed42fd0529eeeb/lean/OAI/Combinatorics/MatroidCounting/CommonBases.lean\#L17689-L17706}{\src{PolySafe}},
\href{https://github.com/openai/math/blob/fd4aeeb2ee4fc729c18d98444fed42fd0529eeeb/lean/OAI/Combinatorics/MatroidCounting/CommonBases.lean\#L22444-L22466}{attribute composition}, and
\href{https://github.com/openai/math/blob/fd4aeeb2ee4fc729c18d98444fed42fd0529eeeb/lean/OAI/Combinatorics/MatroidCounting/CommonBases.lean\#L24280-L24308}{the exponent and estimator attributes}.}

\paragraph{From output laws to successful tapes.}
The final probabilistic interface connects the program to its accuracy guarantee. Let $V$ be the canonical bounded experiment, and let $v$ be the rational output function implemented by the postamble. The certificate asserts
\[
 \E_w\,\varphi(v(w))=\E_b\,\varphi(V(b))
 \quad\text{for every }\varphi:\mathbb Q\to\R,
\]
where both expectations are finite uniform averages. It also records source execution and nonnegativity, including zero output when $Z=0$.\footnote{\href{https://github.com/openai/math/blob/fd4aeeb2ee4fc729c18d98444fed42fd0529eeeb/lean/OAI/Combinatorics/MatroidCounting/CommonBases.lean\#L25962-L25968}{The output certificate}.} The universal observable interface transports the output distribution through changes of enumeration and bank layout.

The machine correspondence reaches actual tape addresses. For $m$ total gathered bank bits, there is a start offset $s$, fixed before quantifying over tapes, such that every tape $\beta:\N\to\{0,1\}$ produces
\[
 v\bigl(\operatorname{physicalWord}(s,m,\beta)\bigr),\qquad
 \operatorname{physicalWord}(s,m,\beta)(i)
   =\beta\bigl(s+3(m-1-i)+1\bigr),\quad i<m.
\]
The reversal comes from pushing bits onto a stack; the stride three reflects the sampling loop's instructions. The proof supplies $s+3m+1\leq B$. Injectivity of these positions shows that a uniform length-$B$ tape induces a uniform length-$m$ word.\footnote{\href{https://github.com/openai/math/blob/fd4aeeb2ee4fc729c18d98444fed42fd0529eeeb/lean/OAI/Combinatorics/MatroidCounting/CommonBases.lean\#L25658-L25687}{Physical tape positions and their uniform law}; \href{https://github.com/openai/math/blob/fd4aeeb2ee4fc729c18d98444fed42fd0529eeeb/lean/OAI/Combinatorics/MatroidCounting/CommonBases.lean\#L26027-L26069}{the machine correspondence}.} Here $m$ counts bank bits, whereas $B$ bounds complete machine execution.

The endpoint quantifies over one fixed program and constants $C,a$ before
the input matroids, oracle, and rational parameters
$0<\varepsilon,\delta<1$. For an exact independence oracle and matroids
with full ground set and supplied common rank $r$, the resource bound is
\[
 B=C\bigl(1+n+\ell+\lceil\varepsilon^{-1}\rceil
       +\operatorname{clog}_2\lceil\delta^{-1}\rceil\bigr)^a,
\]
where $\ell$ is the encoded input length and $\operatorname{clog}_2$ is
the ceiling binary logarithm. Execution on every infinite tape produces
an output within $B$, with at most $B$ oracle calls and bit operations.
For a length-$B$ word, let $\operatorname{out}(w)$ denote this output on
the word's zero extension. Substituting the accuracy indicator in the
observable identity yields
\[
 \#\{w\in\{0,1\}^B:(1-\varepsilon)Z\leq\operatorname{out}(w)
                  \leq(1+\varepsilon)Z\}
 \geq(1-\delta)2^B.
\]
The specification packages this success count with the universal resource and zero-output guarantees.\footnote{\href{https://github.com/openai/math/blob/fd4aeeb2ee4fc729c18d98444fed42fd0529eeeb/lean/OAI/Combinatorics/MatroidCounting/CommonBases.lean\#L110-L132}{The fixed-program and finite-tape specification}.} Verified compilation and probabilistic semantic preservation are established methods \citep{Leroy2009,CakeML2014,CertiCrypt2009,EasyCrypt2011}. Expectation-functional semantics for randomized programs is also established
\citep{AudebaudPaulinMohring2009}. In Lean, SampCert separates sampling
procedures from normalized distribution semantics and proves exact sampler
laws \citep{deMedeirosEtAl2025}. Its \src{SLang} represents a program by an
unnormalized mass function; normalization and equalities between sampler
laws support downstream privacy reasoning. The common-base development
instead carries its bounded random experiment through a costed expression
language to a fixed stack machine. A close compiler precedent is Zar,
which translates probabilistic programs through choice-fix trees and
interaction trees to random-bit samplers, proving equidistribution in
Coq \citep{BagnallStewartBanerjee2023}. Here the comparison turns on the
specification carried through the pipeline: intermediate widths and
numerical attributes, an every-tape resource bound, and the finite count
of accurate tapes. The distinctive evidence is the
concrete integration: the proof carries the counting quantity through oracle code, width bounds, and physical tape coordinates to an exact finite success count. Sections~\ref{sec:tape-example}--\ref{sec:success-budget} isolate the uniformity, law, and error-budget arguments behind this last step.

\subsection{Graph switches: a local table and a fiberwise error direction}
\label{sec:switch-chain}

The graph case studies how two conditional averages interact. Its
switch-chain development connects a four-vertex representation to a
spectral estimate \citep{OAISwitch}. The chain stays put with probability
$1/2$; otherwise it chooses uniformly a four-vertex set and an ordered
pair of distinct perfect matchings on it. It replaces the first matching
by the second when the removed edges are present and the added edges are
absent, and stays put otherwise. Each of the $6\binom n4$ proposals has
probability $1/(12\binom n4)$. For $n\geq4$ and a graphical degree
sequence satisfying $d(v)\leq n-1$, the endpoint gives mixing time at most
$2n^8$ at total-variation threshold $1/4$. When there is more than one
realization, its spectral-gap bound is
\[
 \frac{1}{24n^2\binom n4}.
\]
We follow the disjoint-pair part of the argument, where an exact matrix
count identifies the exceptional direction that the operator estimate
must control.

\paragraph{Reducing a graph cell to a matrix count.}
States are simple graphs on $\operatorname{Fin}n$ realizing $d$. Fix four distinct vertices, grouped into pairs $a=\{i,j\}$ and $b=\{k,l\}$. A coarse cell fixes the outer graph, the internal edge of each pair, and every outside vertex's total incidence into each pair. The remaining data consist of a $2\times2$ Boolean matrix $M$ of cross edges and two sets $U,W$ recording the singleton boundary choices at $i$ and $k$. The source proves that these coordinates determine the graph.\footnote{\href{https://github.com/openai/math/blob/fd4aeeb2ee4fc729c18d98444fed42fd0529eeeb/lean/OAI/Probability/SwitchChain/DisjointEncoding.lean\#L20-L30}{The coarse graph cell}; \href{https://github.com/openai/math/blob/fd4aeeb2ee4fc729c18d98444fed42fd0529eeeb/lean/OAI/Probability/SwitchChain/DisjointCoordinateFeasibility.lean\#L22-L108}{degree feasibility in coordinates}.}

Once the cell data are fixed, the four degree equations reduce to the
cross-edge total and two margin equations:
\[
 \operatorname{row}_0(M)+|U|=c_L,\qquad
 \operatorname{col}_0(M)+|W|=c_R,
\]
where $c_L,c_R$ are determined by the original graph. Forward and inverse
constructions identify the graph cell with boundary choices and a
compatible matrix. Under this equivalence, a uniform graph law weights
each pair of boundary choices by its number of compatible matrices.
Preserving these multiplicities is what makes the coordinate model usable
for conditional averaging. The matrix count describes the remaining
dependence between the left and right choices.

There are only sixteen Boolean matrices. Let $c_s(r,t)$ count those with total $s$, first-row sum $r$, and first-column sum $t$. For $s=2$, finite decision proves the table
\[
 \bigl(c_2(r,t)\bigr)_{r,t=0}^2
   =\begin{pmatrix}0&1&0\\1&2&1\\0&1&0\end{pmatrix},
 \qquad c_2(r,t)=2-p(r)-p(t),
\]
where $p(r)=\mathbf1_{r\ne1}$.\footnote{\href{https://github.com/openai/math/blob/fd4aeeb2ee4fc729c18d98444fed42fd0529eeeb/lean/OAI/Probability/SwitchChain/CrossMatrix.lean\#L22-L25}{The finite total-two count}.}
The other totals give tables that factor into independent margin weights.
For total two, centering leaves dependence through one Boolean flag on
each side. This reduces the correlated component to a specified local
direction, which the subsequent operator argument can isolate.

A short calculation illustrates the exceptional correlation. Among the six
uniformly chosen total-two matrices, each off-center margin flag has mean
$1/3$, while their product is always zero. Their covariance is therefore
$-1/9$. The sign shows why a direct nonnegative-correlation argument is
unavailable. Actual graph cells also carry boundary multiplicities, so the
proof transports the flag description to those weighted fibers before
estimating its contribution.

\paragraph{Paying for the exceptional correlation.}
In the finite coordinate model, composing the two conditional averages
gives $A+Bp(r)$ on each fixed cell. The coefficients $A,B$ depend on
that cell and the averaged function, and $B=0$ unless $s=2$. The graph-cell equivalence transports this formula back to actual graph states. This transport includes a proof that corresponding fibers have equal cardinalities and corresponding sums, and an identification of the flag: $p(r)$ says that $i$ has equal adjacency bits to $k$ and $l$, meaning both cross edges are present or both are absent.\footnote{\href{https://github.com/openai/math/blob/fd4aeeb2ee4fc729c18d98444fed42fd0529eeeb/lean/OAI/Probability/SwitchChain/CrossCellAffineComposition.lean\#L64-L87}{The affine conditional image}; \href{https://github.com/openai/math/blob/fd4aeeb2ee4fc729c18d98444fed42fd0529eeeb/lean/OAI/Probability/SwitchChain/DisjointGraphAffine.lean\#L68-L90}{its graph interpretation}.}

Work in the real function space on the finite graph state space $\Omega$,
with counting inner product $\langle f,g\rangle=\sum_{G\in\Omega}f(G)g(G)$.
Fiber averages themselves are divided by the fiber cardinality.\footnote{See
\href{https://github.com/openai/math/blob/fd4aeeb2ee4fc729c18d98444fed42fd0529eeeb/lean/OAI/Probability/SwitchChain/OrthogonalPairResampling.lean\#L191-L193}{\src{OrthogonalPairResampling.lean}, lines 191--193}, and
\href{https://github.com/openai/math/blob/fd4aeeb2ee4fc729c18d98444fed42fd0529eeeb/lean/OAI/Probability/SwitchChain/OrthogonalPairResampling.lean\#L231-L232}{lines 231--232}.}
Write $E_a$ for averaging over a
pair-resampling fiber. The transported formula expresses the relevant conditional image as a fiber-constant base plus a fiber-constant weight times the equality flag; the weight is supported where the cross total is two. For this purpose a fiber is active when $k$ and $l$ each have exactly one neighbor in $a$. Subtracting $E_a$ leaves a multiple of the centered flag on each active fiber. These functions form an error subspace whose orthogonal projection is denoted $R_{ab}$. The description is at most one dimensional on each fiber; the complete graph-space subspace can have many independent fiber components.

The flag representation now supplies the hypotheses of an abstract
projection inequality. Conditional averaging is symmetric and idempotent,
and the error projection fixes the correlated fluctuation. Decomposition
into the projections' ranges and orthogonal complements turns these facts
into an interaction bound. If $G_{ab}$ denotes the switch-gain projection,
the graph instance gives
\[
 2\norm{G_{ab}f}^2-\norm{R_{ab}f}^2-\norm{R_{ba}f}^2
 \leq2\langle(I-E_a)f,(I-E_b)f\rangle.
\]
To sum this inequality, the proof needs a budget for the error terms.
Fix a graph
$G$ and a pair $a=\{i,j\}$. Its singleton-neighbor set $V_G$ consists of
outside vertices adjacent to exactly one endpoint of $a$. The vertices
assigned to $i$ form a subset of prescribed cardinality $q$. An equivalence
therefore identifies this assignment space with the slice
$\{S\subseteq V_G:|S|=q\}$. Under its induced linear isometry, the
centered equality flag becomes the slice's \src{raw} vector, while a
graph switch becomes a transposition of two labels.

The slice estimate is itself an operator comparison. Normalized flag
vectors have squared norm $1/2$ whenever nonzero. Their synthesis operator
$T$ satisfies $TT^*\leq L_{\mathrm{swap}}$ in quadratic-form order.
The proof decomposes coefficient data into a zero-sum linear part and a
symmetric part with zero rows and diagonal. The Gram operator is scalar
on these parts; its scalars are bounded by the corresponding swap-generator
scalars. Boundary slices have zero flag vectors and are treated separately.
The decomposition controls the overlap between flag directions. In
projection form it gives
\[
 \frac12\sum_{k,l}\norm{\Pi_{kl}h}^2
 \leq\frac18\sum_{k,l}\norm{h-h\circ\tau_{kl}}^2,
\]
where $\Pi_{kl}$ projects onto the centered flag line and $\tau_{kl}$
swaps the two labels.\footnote{See
\href{https://github.com/openai/math/blob/fd4aeeb2ee4fc729c18d98444fed42fd0529eeeb/lean/OAI/Probability/SwitchChain/SliceOperatorComparison.lean\#L96-L156}{the two-block operator comparison}, and
\href{https://github.com/openai/math/blob/fd4aeeb2ee4fc729c18d98444fed42fd0529eeeb/lean/OAI/Probability/SwitchChain/SliceReindexedComparison.lean\#L49-L79}{isometric reindexing and the swap-sum estimate}.}

Two normalizations connect this estimate to the graph budget. Multiplying
it by two gives a coefficient $1/4$ on the squared differences, exactly
the squared norm of the switch-gain projection $(I-S)/2$ for the switch
involution $S$. Next, the restrictions $h_G$ to assignment spaces satisfy
the exact identity
\[
 \sum_{G\in\Omega}
 \frac{\norm{h_G}^2}{\binom{|V_G|}{q_G}}=\norm h^2.
\]
Each fiber is represented once by each of its graph states; division by its
cardinality cancels that repetition. Applying this identity to both sides
of the assignment budget lifts the local inequality to the whole graph
space. Finally, ordered partner labels $(k,l)$ count each unordered pair
twice, so both error and gain sums acquire the same factor $1/2$.
The resulting budget bounds the sum of squared error norms by the sum of
squared gain norms.\footnote{The restriction identity and its use are in
\href{https://github.com/openai/math/blob/fd4aeeb2ee4fc729c18d98444fed42fd0529eeeb/lean/OAI/Probability/SwitchChain/PairEqualityBudget.lean\#L151-L270}{\src{PairEqualityBudget.lean}, lines 151--270}; the ordered/unordered conversion is in
\href{https://github.com/openai/math/blob/fd4aeeb2ee4fc729c18d98444fed42fd0529eeeb/lean/OAI/Probability/SwitchChain/PairEqualityBudgetUnordered.lean\#L23-L54}{\src{PairEqualityBudgetUnordered.lean}, lines 23--54}.}
Summing the compensated inequality over ordered disjoint two-element vertex sets $a,b$ then gives
\[
 \sum_{a,b:\,a\cap b=\varnothing}
   \langle(I-E_a)f,(I-E_b)f\rangle\geq0.
\]
The compensated local bounds therefore give a nonnegative aggregate
interaction even though individual fibers can exhibit negative
correlation.\footnote{\href{https://github.com/openai/math/blob/fd4aeeb2ee4fc729c18d98444fed42fd0529eeeb/lean/OAI/Probability/SwitchChain/PairActualAffine.lean\#L81-L101}{The concrete interaction inequality}; \href{https://github.com/openai/math/blob/fd4aeeb2ee4fc729c18d98444fed42fd0529eeeb/lean/OAI/Probability/SwitchChain/PairDisjointAssembly.lean\#L45-L77}{assembly with the projection budget}.}

Triple-generator and local comparison estimates combine with the
disjoint sum to give a generator-square bound, from which the stated gap
and mixing estimate follow. Finite computation and transport through
equivalent representations are established proof methods
\citep{Harrison1995,CohenMahboubi2012}. The informative integration here
is the route from a sixteen-object count to a global estimate: identify
the equality flag, transport its centered direction with the correct
multiplicities, and pay for its aggregate energy through switch gains.
